% -----------------------------------------------------------
%  Recopilación de Información
% -----------------------------------------------------------

Este anexo describe las técnicas de recopilación de información empleadas durante las etapas de análisis y diseño del proyecto DOMU. La información recopilada alimentó la definición de requerimientos funcionales, la especificación de casos de uso y las decisiones de diseño del sistema.

% === PLACEHOLDER: Reemplazar con imagen generada ===
% PROMPT para generador de imágenes (Gemini/GPT/DALL-E):
%
% "Genera un DIAGRAMA DE FLUJO DE INFORMACIÓN mostrando las 4 fuentes de recopilación de datos
% para un proyecto de software.
%
% NOTACIÓN:
% - FUENTES: 4 rectángulos redondeados con borde sólido, cada uno con un título en negrita y
%   debajo una lista con viñetas (bullet points) de los sub-ítems.
% - ELEMENTO CENTRAL: un rectángulo más grande con borde doble (o más grueso), en el centro,
%   etiquetado 'Requerimientos Funcionales (20 RFs) + Casos de Uso (27 CUs)'.
% - FLECHAS: flechas gruesas con punta, desde cada fuente hacia el centro, indicando que la
%   información fluye de las fuentes a los requerimientos.
% - ÉNFASIS: la caja 'Análisis de Competencia' debe tener borde más grueso y un asterisco (*)
%   o etiqueta 'Fuente principal' para destacarla visualmente respecto a las otras 3.
%
% DISPOSICIÓN (las 4 fuentes rodean al centro):
% ARRIBA-IZQUIERDA: 'Revisión de Documentos Externos'
%   • Ley 21.442 Copropiedad
%   • Ley 21.719 Datos Personales
%   • Ley 19.799 Docs Electrónicos
%   • Estadísticas Subtel
%
% ARRIBA-DERECHA: 'Análisis de la Competencia' ★ FUENTE PRINCIPAL (borde grueso doble)
%   • ComunidadFeliz
%   • Edifito
%   • Edipro
%   • GastoComunChile
%
% ABAJO-IZQUIERDA: 'Reuniones con Profesora Guía'
%   • Validación de alcance
%   • Retroalimentación de diseño
%   • Revisión de avances
%
% ABAJO-DERECHA: 'Observación de Procesos en Condominios'
%   • Libro de visitas manual
%   • Pizarras de avisos
%   • Gestión de encomiendas en cuaderno
%   • Reservas en papel
%
% ESTILO: Fondo blanco, bordes negros, relleno blanco o gris muy claro. Sin gradientes.
% Tipografía sans-serif. Profesional para tesis universitaria. Etiquetas en español.
% Formato: PNG 300dpi."
\begin{figure}[H]
  \centering
  \fbox{\parbox{0.85\textwidth}{\centering\vspace{3cm}
  \textit{[Insertar diagrama: Fuentes de recopilación de información y su relación con los RFs y CUs del proyecto]}
  \vspace{3cm}}}
  \caption{Técnicas de recopilación de información y su aporte a los requerimientos del sistema.}
  \label{fig:fuentes-recopilacion}
\end{figure}

\section{Revisión de Documentos Externos}

Se revisaron documentos normativos y estadísticos relevantes para el dominio de la administración de comunidades residenciales. La Tabla~\ref{tab:documentos-externos} detalla las fuentes consultadas y la información útil extraída de cada una.

\begin{longtable}{|p{0.28\textwidth}|p{0.18\textwidth}|p{0.46\textwidth}|}
\caption{Documentos externos revisados}\label{tab:documentos-externos}\\
\hline
\textbf{Documento} & \textbf{Fuente / Fecha} & \textbf{Información útil extraída} \\
\hline
\endfirsthead
\hline
\textbf{Documento} & \textbf{Fuente / Fecha} & \textbf{Información útil extraída} \\
\hline
\endhead
\hline
\endfoot

Ley 21.442 de Copropiedad Inmobiliaria y su Reglamento (D.S. N°\,1 de 2025) &
Biblioteca del Congreso Nacional / 2022--2025 &
Obligaciones del administrador (rendición de cuentas, fondo de reserva, sesiones de copropietarios), derechos de los residentes, quórum de votación. Estos elementos fundamentaron los RF de finanzas (RF\_07, RF\_08), votaciones (RF\_10) y RBAC (RF\_16). \\
\hline

Ley 21.719 de Protección de Datos Personales &
Biblioteca del Congreso Nacional / 2024 &
Requisitos de consentimiento para almacenamiento de datos personales, derecho de acceso y eliminación. Influyó en las decisiones de diseño del módulo de registro y gestión de perfiles. \\
\hline

Ley 19.799 de Documentos Electrónicos &
Biblioteca del Congreso Nacional / 2002 &
Validez jurídica de documentos electrónicos y firma electrónica. Relevante para la generación de boletas y estados de cuenta en formato PDF (RF\_07). \\
\hline

Estadísticas de penetración de Internet en Chile &
Subtel / 2024 &
Datos de conectividad residencial (cobertura de banda ancha fija y móvil en zonas urbanas). Respaldó la decisión de implementar una aplicación web progresiva (PWA) como canal principal de acceso. \\

\end{longtable}

\section{Análisis de la Competencia}

El análisis de la competencia constituyó la \textbf{fuente principal} de recopilación de información para el diseño funcional de DOMU. Se evaluaron cuatro plataformas chilenas de administración de comunidades, con el objetivo de identificar funcionalidades estándar, brechas de mercado y oportunidades de diferenciación.

\subsection{Plataformas evaluadas}

\begin{enumerate}
    \item \textbf{ComunidadFeliz} --- Plataforma líder en Chile con cobertura integral: gastos comunes, reservas, comunicación, control de acceso. Modelo SaaS con planes desde UF~0.10 por unidad/mes.

    \item \textbf{Edifito} --- Solución orientada a la gestión financiera y comunicación entre vecinos. Incluye módulo de pagos en línea y publicación de avisos.

    \item \textbf{Edipro} --- Plataforma enfocada en la administración profesional de edificios. Ofrece gestión de gastos comunes, cobranza y reportes financieros.

    \item \textbf{GastoComunChile} --- Solución especializada en el cálculo y distribución de gastos comunes, con énfasis en el cumplimiento de la Ley de Copropiedad.
\end{enumerate}

\subsection{Criterios de evaluación}

Las plataformas se evaluaron según los siguientes criterios:
\begin{itemize}
    \item \textbf{Funcionalidades}: Módulos disponibles (finanzas, reservas, comunicación, control de acceso, proveedores, votaciones).
    \item \textbf{Modelo de negocio}: Estructura de precios, costos por unidad habitacional, planes disponibles.
    \item \textbf{Experiencia de usuario}: Calidad de la interfaz, disponibilidad móvil, facilidad de uso.
    \item \textbf{Limitaciones}: Funcionalidades ausentes, restricciones técnicas, dependencia de hardware especializado.
\end{itemize}

\subsection{Resultados y brechas identificadas}

La tabla comparativa detallada se presenta en el Capítulo~\ref{cap:introduccion}, Sección~1.4. Las principales brechas identificadas que motivaron los requerimientos funcionales de DOMU fueron:

\begin{itemize}
    \item \textbf{Control de acceso con QR de bajo costo:} Las soluciones existentes requieren hardware propietario (tótem, tablets dedicadas). DOMU propone un modelo con pistola lectora QR estándar (modo HID), reduciendo significativamente el costo de implementación (RF\_01, RF\_03).

    \item \textbf{Chat en tiempo real entre vecinos:} La mayoría de las plataformas limitan la comunicación a foros o tablones de avisos. DOMU incorpora chat WebSocket bidireccional (RF\_17).

    \item \textbf{Marketplace comunitario:} Funcionalidad no ofrecida por las plataformas evaluadas. Permite la compraventa de artículos entre vecinos de la misma comunidad (RF\_18).

    \item \textbf{Portal de proveedores:} Las plataformas analizadas gestionan proveedores internamente sin otorgarles acceso directo. DOMU ofrece un portal dedicado con flujo de órdenes y cotizaciones (RF\_13).

    \item \textbf{Biblioteca de documentos:} Centralización de reglamentos, actas y documentos comunitarios con almacenamiento en la nube, funcionalidad ausente o limitada en las alternativas evaluadas (RF\_19).

    \item \textbf{Modelo de código abierto:} A diferencia de las soluciones comerciales evaluadas, DOMU se concibe como un proyecto de código abierto, permitiendo su personalización y despliegue sin costos de licenciamiento.
\end{itemize}

\section{Reuniones con Profesora Guía}

Se realizaron reuniones periódicas con la profesora guía a lo largo del proyecto, con los siguientes objetivos:

\begin{itemize}
    \item \textbf{Validación de alcance:} Revisión de los requerimientos funcionales y priorización de casos de uso según la metodología MoSCoW.
    \item \textbf{Retroalimentación sobre diseño:} Evaluación de decisiones arquitectónicas (elección de Javalin sobre Spring Boot, estrategia de multi-tenancy, modelo de datos financiero).
    \item \textbf{Revisión de avances:} Seguimiento del progreso por fases, con énfasis en la decisión de diferir funcionalidades a v2.0 (semana~8).
    \item \textbf{Revisión del informe:} Correcciones de estructura, contenido y formato del documento final según el template UBB.
\end{itemize}

Las reuniones se realizaron de forma presencial y remota, con una frecuencia aproximada de una sesión cada dos semanas. La retroalimentación recibida influyó directamente en la re-priorización de funcionalidades y en la estructura del informe.

\section{Observación de Procesos en Condominios}

Se realizó una observación informal de los procesos manuales vigentes en edificios de la ciudad de Chillán, con el fin de identificar puntos de dolor y oportunidades de digitalización. Las observaciones se centraron en:

\begin{itemize}
    \item \textbf{Control de acceso:} Libro de visitas en papel, donde el conserje registra manualmente el nombre, RUT y hora de ingreso de cada visitante. Se identificó que este proceso demora entre 2 y 5~minutos por visita y es propenso a errores de registro.

    \item \textbf{Comunicación entre vecinos:} Pizarras de avisos en áreas comunes y grupos de WhatsApp como canales principales de comunicación. Se observó fragmentación de la información y dificultad para acceder a avisos históricos.

    \item \textbf{Gestión de encomiendas:} Registro informal en cuadernos o sin registro, con notificación verbal o mediante nota bajo la puerta. Se identificó pérdida de trazabilidad y demoras en la entrega.

    \item \textbf{Reserva de áreas comunes:} Cuadernos de reserva en conserjería, con conflictos frecuentes por doble reserva y falta de visibilidad del calendario.

    \item \textbf{Gastos comunes:} Distribución de boletas en papel o por correo electrónico manual, con pagos presenciales o transferencia bancaria sin confirmación automática.
\end{itemize}

Estas observaciones alimentaron directamente los diagramas BPMN del Capítulo~\ref{cap:introduccion} (procesos AS-IS) y la definición de los flujos TO-BE que implementa DOMU.